% Owner: classification_e. Root and boundary_d full-source referee PASS.
% Uses the independently refereed abstract purity lemma E.2.11.
\begin{proposition}[E.2.12 (Proved): no additional central critical values]
In the smooth-specialization setting of E.2.10, let
\[
 g:\mathcal X\longrightarrow\mathbb P^1\times\Delta
\]
be the extended intrinsic fibration, and let
$c_1(s),c_2(s),c_0(s)$ be the three labelled critical-value
sections of E.2.7.  The critical values of $g_0$ are exactly
the distinct points among $c_1(0),c_2(0),c_0(0)$.
In particular, no additional critical value occurs only
on the central target curve.

If the normalizing projective transformations $A_s$ of
E.2.7 have a degenerate limit, the central regular fibres
are isomorphic complex two-tori and $g_0$ has at most two
critical values.  If the constant image of the limiting
rank-one matrix is $0$ or $1$, this torus is the
corresponding finite-root abelian surface.  If that image
is an ordinary source point, it is the ordinary torus at
that point and need not be an abelian surface.  The cusp
cannot be the constant image.
\end{proposition}
\begin{proof}
E.2.7 gives $0<|z(0)|<R_{k,l}$.  A local logarithm
therefore extends the normalized parameter $c(s)$
holomorphically to the disc.  The critical image is closed
and contains the three punctured sections, so it contains
their three limiting points.  Suppose that a different
central critical value $b_0=(v_0,0)$ exists.  Take a small
bidisc $B$ about $b_0$ which avoids all three section
closures and all other central critical values.  Put
$U=B\setminus\{b_0\}$.  The restricted family over $U$
is a proper smooth family of complex two-tori.  We verify
that its Jacobian extends to a smooth holomorphic torus
group family over $B$, so that E.2.11 applies.

The normalized source coordinate is $t=A_s(v)$.
Represent $A_s$ by a matrix with holomorphic entries,
after clearing poles and removing a common power of
$s$, so its central matrix is nonzero.  If that matrix
is invertible, $t$ extends holomorphically near $b_0$ and
its value there is ordinary: a marked value would put
$b_0$ on the corresponding section closure.  The
ordinary period model with parameter $c(s)$ then gives
the desired smooth Jacobian over $B$.

Otherwise the central matrix has rank one.  Denote its
constant image by $t_*$ and its exceptional target point
by $v_*$.  For each marked source value different from
$t_*$, its preimage tends to $v_*$.  At least two of the
three marked values are different from $t_*$, so $v_*$
belongs to at least two labelled section closures.
Consequently $v_0\ne v_*$.  After shrinking $B$, the
function $t(s,v)$ tends uniformly to $t_*$ there.
For an ordinary $t_*$ the ordinary period model again
extends directly.

Suppose instead that $t_*=0$ or $1$, and use a local
coordinate $\xi$ at that marked point.  The divisor of
$\xi(t(s,v))$ on $B$ is supported on $s=0$, since $B$
avoids the corresponding critical-value graph.
Consequently
\[
 \xi(t(s,v))=s^n u(s,v),\qquad n>0,
\]
where $u$ is a nowhere-zero holomorphic function.
This also explains why no zero of the leading
coefficient may be overlooked: such a zero would put
the point on the labelled graph closure.
Choose a nearby $v_1$ for which $(v_1,0)$ is regular.
The actual family over the full $s$-disc at $v_1$ is
smooth, so its integral degree-one monodromy is the
identity.  The parameter $c(s)$ has no winding.  By
the marked comparison of E.2.7, this monodromy is
$A_1^n$ or $A_2^n$, respectively.  Their orders are
three and four, so $e\mid n$, where $e=3$ or $4$.
A holomorphic $e$th root of $u$ exists on the bidisc,
and
\[
 r(s,v)=s^{n/e}u(s,v)^{1/e}
\]
is an actual holomorphic root coordinate on $B$.
Pull back the smooth finite-root torus period model
by $(c(s),r(s,v))$.  This produces the required
smooth torus group family $G\to B$.  The larger
multiplicities $3k,4l$ do not change this Jacobian
argument: their additional deck transformations
are translations, and the linear Jacobian model
is the original order-three or order-four root model.

For completeness $t_*=\infty$ is impossible.  In the
cusp coordinate the same divisor argument gives
$q=s^n u$ with $n>0$.  The smooth family at $v_1$
would have monodromy $A_0^n$.  But
$A_0=I+N$, $N^2=0$, and
$N\widehat u=\widehat w\ne0$, so $A_0^n\ne I$.

It remains to identify $G|_U$ with the actual Jacobian,
rather than merely comparing individual tori.
The punctured bidisc $U$ is simply connected.  Its
integral first-homology local system therefore has
a global marking.  Over $s\ne0$ that marking agrees
with the reference period marking, after one constant
integral change of basis; the transition comparisons
of E.1.3 have identity linear homology action at
fixed marked parameter.  Both holomorphic period
subbundles extend on $U$, and they agree on its dense
open subset $s\ne0$.  They agree everywhere on $U$.
Equivalently, the two quotients of their holomorphic
Lie bundles by the same marked period lattice are
isomorphic torus group families on $U$.  This is a
Jacobian identification and makes no assertion that
the actual affine torsor has a section on $B$.

The total space is smooth and $g$ is proper and flat.
The only possible nonsmooth fibre over $B$ lies above
$b_0$, and its codimension in the total space is two.
All hypotheses of E.2.11 now hold.  That lemma makes
$g$ smooth also above $b_0$, a contradiction.  The
central critical-value set is therefore precisely
the set of labelled limits.

For the final assertions, in the rank-one case every
ordinary central target point different from $v_*$
has the same limiting marked torus: it is obtained
from $(c(0),t_*)$, or from the finite-root value
$(c(0),r=0)$.  Equality of the extended period
subbundles proves this assertion without assuming
separatedness for threefold completions.
All marked values except possibly $t_*$ have their
preimages tending to $v_*$.  Hence there are at most
two central critical values; if $t_*$ is ordinary,
there is exactly one.
At $t_*=0$ or $1$, the finite linear automorphism of
the root torus has complex eigenvalues $1$ and a
primitive third or fourth root of unity.  The
rational idempotents for these two distinct
cyclotomic factors split that torus, up to isogeny,
as a product of two elliptic curves.  It is therefore
an abelian surface.  This argument makes no
projectivity assertion for an ordinary limiting
torus, and supplies no polarization on the whole
central threefold.
\end{proof}

\begin{remark}[Scope of E.2.11--E.2.12]
Purity here is obtained from the explicit good
Jacobian near a hypothetical extra critical value.
Smooth total space and flatness alone were not used
as a general purity theorem.  These results do not
give uniqueness of an interior smooth completion,
do not exclude the remaining isotrivial central
fibrations, and do not prove persistence of the
canonical pencil on an arbitrary new deformation
branch.
\end{remark}
